\documentclass[10pt,a4paper]{amsart}
\usepackage[margin=19mm]{geometry}
\usepackage{amsmath,amssymb,mathtools}
\usepackage[scr=rsfs,cal=euler]{mathalfa}
\usepackage{xfrac}
\usepackage[unicode,hidelinks,bookmarks=false]{hyperref}
\newcommand{\ie}{i.e.\ }
\newcommand{\cf}{cf.\ }
\newcommand{\resp}{respectively\ }
\newcommand{\Subsection}{\subsection}
\providecommand{\nobreakdash}{\nobreak\hbox{-}}
\setlength{\emergencystretch}{2em}
\multlinegap0pt
\usepackage{float}
\usepackage{amssymb}
\newtheorem{lemma}{Lemma}
\newtheorem{definition}{Definition}
\newtheorem{theorem}{Theorem}
\title{Spectral Factorization and Hypergeometric Representations of the Alexander Polynomials of $Th(4,2n+1)$}
\author{Suman Saurabh}
\date{Source: arXiv:2606.11301v1 (CC BY 4.0)}
\begin{document}
\maketitle

\noindent\emph{Source and licence.} Spectral Factorization and Hypergeometric Representations of the Alexander Polynomials of $Th(4,2n+1)$. Version arXiv:2606.11301v1 (CC BY 4.0), distributed by arXiv under the Creative Commons Attribution 4.0 International licence, http://creativecommons.org/licenses/by/4.0/, as recorded in the arXiv licence metadata for this exact version and shown on the abstract page https://arxiv.org/abs/2606.11301v1. Full licence notice: \texttt{licenses/arxiv-2606.11301-CC-BY-4.0.txt}.\par\smallskip

\noindent\textit{Editorial scope.} Counted unit: the article's theorem thm:spectral\_fact (source0.tex lines 301--306). The article's complete original proof of it is delivered with it (source0.tex lines 308--347, one \texttt{proof} environment of 40 lines). No mathematics is written, repaired or re-derived: the statement and the proof are the article's own lines, in the article's own notation, and no retained line is substituted or re-flowed.  Retained uncounted as the source-local context the proof needs: lines 206--211; lines 217--220; lines 236--242; lines 251--256; lines 266--271; lines 285--291.  Nothing was omitted from any retained span, so the package carries no omission record.  No retained line contains a citation, so the wrapper carries no bibliography and no bibliography entry.  No retained line contains a figure environment, an image-inclusion command or drawing code, so no image is delivered and the package declares no asset dependency.  Editorial formatting only, in addition: an AMS \texttt{amsart} preamble that restates the article's own declarations and macros used by the retained lines, namely the environments \texttt{lemma}, \texttt{definition}, \texttt{theorem} on separate counters, together with the standard packages those lines need.  The retained environments print in this document as Lemma 1, Definition 1, Theorem 1 in order of appearance, whereas the article prints the same items with the numbering of its own full document.  The complete native source of the article as distributed in the arXiv e-print for this exact version is bundled unchanged as \texttt{sources/202538/source0.tex}.
\begin{lemma}[Determinant Factorization]\label{lem:det_fact}
Let $q = 2n+1$. Over the $q$-th roots of unity $\zeta^j = \exp\left(\frac{2\pi i j}{q}\right)$, the Burau determinant expands as:
\begin{equation}
\det(I - M^q) = \prod_{j=0}^{2n} \det(I - \zeta^{-j} M).
\end{equation}
\end{lemma}

\begin{lemma}[Trivial Root Evaluation]\label{lem:trivial_root}
The trivial root $j=0$ evaluates to:
\[ \det(I - M(t)) = t^{-1}(1 + t + t^2 + t^3). \]
\end{lemma}

\begin{definition}[Normalized Polynomial Representation]\label{def:recip_poly}
The Burau determinant identity defines the Alexander polynomial up to a unit $\pm t^k$. We fix this ambiguity by defining the normalized ordinary polynomial $A_{2n+1}(z)$ to be the polynomial resulting from the unnormalized determinant expansion evaluated at $t=-z$ and scaled by $z^{2n+1}$:
\begin{equation}
A_{2n+1}(z) = -z^{2n+1} \frac{\det(I - M(-z)^{2n+1})}{1-z+z^2-z^3}.
\end{equation}
This selects a unique representative of the Alexander equivalence class. We define its corresponding symmetric Laurent representative as $\nabla_{2n+1}(z) = z^{-3n} A_{2n+1}(z)$.
\end{definition}

\begin{lemma}[Reciprocity Relation]\label{lem:reciprocity}
The characteristic polynomial satisfies the reciprocity identity:
\[
\chi_M(x^{-1},-z) = -zx^{-3}\chi_M(x,-z^{-1}).
\]
\end{lemma}

\begin{lemma}[Resultant Polynomial]\label{lem:resultant}
Let $v = x + x^{-1}$ denote the reciprocal spectral parameter. Clearing the Laurent denominator by defining $\widetilde{\chi}_M(x, w) = w \chi_M(x, -w)$, the resultant polynomial over the constraint $w^2-uw+1=0$ evaluates to:
\begin{equation}
R(x, u) = \operatorname{Res}_w\big(\widetilde{\chi}_M(x, w), w^2-uw+1\big) = -x^3 (u-v)(u-v+2)^2.
\end{equation}
\end{lemma}

\begin{lemma}[Chebyshev Roots]\label{lem:cheb_roots}
The characteristic Chebyshev combination 
\begin{equation}
F_n(u) = U_n\left(\frac{u}{2}\right) + U_{n-1}\left(\frac{u}{2}\right)
\end{equation}
is monic of degree $n$ and factors as $F_n(u) = \prod_{r=1}^n (u - v_r)$.
\end{lemma}

\begin{theorem}[Spectral Factorization]\label{thm:spectral_fact}
The symmetric Laurent representative admits a factorization in $\mathbb{Z}[u]$. Defining $B_{2n+1}(u) = \nabla_{2n+1}(z)$, the polynomial evaluates to:
\begin{equation}
B_{2n+1}(u) = F_n(u)F_n(u+2)^2.
\end{equation}
\end{theorem}

\begin{proof}
Evaluating the unnormalized determinant identity defined in Definition \ref{def:recip_poly} yields:
\begin{equation*}
A_{2n+1}(z) = \frac{-z^{2n+1} \det(I - M(-z)^{2n+1})}{1-z+z^2-z^3}.
\end{equation*}
By Lemma~\ref{lem:trivial_root}, the trivial root corresponding to $j=0$ evaluates to $\det(I-M(-z)) = -z^{-1}(1-z+z^2-z^3)$. Canceling this from the product expansion in Lemma~\ref{lem:det_fact} isolates the contribution of the non-trivial roots:
\begin{equation*}
A_{2n+1}(z) = z^{2n} \prod_{j=1}^{2n} \det(I - \zeta^{-j}M(-z)).
\end{equation*}
By the definition of the characteristic polynomial, $\det(cI - M) = c^3\chi_M(c^{-1})$. Setting $c = \zeta^{-j}$, the determinant evaluates to $\det(\zeta^{-j}I - M(-z)) = \zeta^{-3j}\chi_M(\zeta^j, -z)$. Since $\det(I - \zeta^{-j}M) = \det(\zeta^{-j}(\zeta^j I - M))$, this yields:
\begin{equation*}
\det(I - \zeta^{-j}M(-z)) = \zeta^{-3j}\chi_M(\zeta^j, -z).
\end{equation*}
The summation of the indices in the prefactor evaluates to $\sum_{j=1}^{2n} -3j = -3n(2n+1)$. Since $\zeta$ is a $(2n+1)$-th root of unity, $\zeta^{-3n(2n+1)} = (\zeta^{2n+1})^{-3n} = 1$. Consequently, the product of the non-trivial roots reduces exactly to the product of the characteristic polynomials:
\begin{equation*}
A_{2n+1}(z) = z^{2n} \prod_{j=1}^{2n} \chi_M(\zeta^j, -z).
\end{equation*}
We partition the $2n$ roots of unity into $n$ reciprocal pairs $(\zeta^r, \zeta^{-r})$ for $1 \le r \le n$. The sequence decomposes as:
\begin{equation*}
A_{2n+1}(z) = z^{2n} \prod_{r=1}^{n} \Big( \chi_M(\zeta^r, -z) \chi_M(\zeta^{-r}, -z) \Big).
\end{equation*}
Applying the reciprocity identity (Lemma~\ref{lem:reciprocity}) to the second factor, we substitute \newline
$\chi_M(\zeta^{-r}, -z) = -z\zeta^{-3r}\chi_M(\zeta^r, -z^{-1})$. The product for a single pair index $r$ becomes:
\begin{equation*}
\chi_M(\zeta^r, -z) \Big( -z\zeta^{-3r}\chi_M(\zeta^r, -z^{-1}) \Big) = -z\zeta^{-3r} \Big[ \chi_M(\zeta^r, -z)\chi_M(\zeta^r, -z^{-1}) \Big].
\end{equation*}
By Lemma~\ref{lem:resultant}, the bracketed term evaluates to the resultant $R(x, u)$ evaluated at $x = \zeta^r$ and $u = z + z^{-1}$. Substituting $R(\zeta^r, u) = -(\zeta^r)^3 (u - v_r)(u - v_r + 2)^2$, where $v_r = \zeta^r + \zeta^{-r}$, yields:
\begin{equation*}
-z\zeta^{-3r} \Big[ -\zeta^{3r} (u - v_r)(u - v_r + 2)^2 \Big] = z(u - v_r)(u - v_r + 2)^2.
\end{equation*}
Taking the product over all $n$ reciprocal pairs incorporates the $z^{2n}$ prefactor:
\begin{equation*}
A_{2n+1}(z) = z^{2n} \prod_{r=1}^n z(u - v_r)(u - v_r + 2)^2 = z^{3n} \prod_{r=1}^n (u - v_r)(u - v_r + 2)^2.
\end{equation*}
Dividing both sides by $z^{3n}$ produces the symmetric Laurent representative $\nabla_{2n+1}(z)$:
\begin{equation*}
\nabla_{2n+1}(z) = z^{-3n} A_{2n+1}(z) = \prod_{r=1}^n (u - v_r)(u - v_r + 2)^2.
\end{equation*}
The right-hand side is a polynomial in $u = z+z^{-1}$. Let $P(u) = \prod_{r=1}^n (u - v_r)(u - v_r + 2)^2$. Because $(u - v_r)$ is monic of degree 1 and $(u - v_r + 2)^2$ is monic of degree 2, their product over $n$ iterations ensures $P(u)$ is identically monic of degree $3n$. By Lemma~\ref{lem:cheb_roots}, the Chebyshev polynomial $F_n(u)$ is monic of degree $n$ with exact roots $v_r$. Thus, the polynomial $F_n(u)F_n(u+2)^2$ is monic of degree $3n$ and shares the root factorization of $P(u)$, establishing $B_{2n+1}(u) = F_n(u)F_n(u+2)^2$.
\end{proof}

\end{document}
